\chapter{Paracentesis}
\index{Paracentesis}

\medskip
\noindent\textit{Educational information; seek professional advice for individual care.}

\section*{What it is}
Paracentesis is a procedure in which a needle or drain is passed through the abdominal wall to remove ascites, the build-up of fluid in the abdomen. It may relieve pressure or breathlessness, or allow fluid to be tested to find the cause of the ascites.

\section*{Before the procedure}
The team reviews the reason for the procedure, medicines including anticoagulants, allergies, pregnancy where relevant, infection risk, and blood tests. Ultrasound may help identify a safe site. Do not stop a prescribed medicine unless the treating team gives instructions.

\section*{What to expect}
The skin is cleaned and numbed with local anaesthetic. A sample or a larger volume of fluid may be removed, sometimes over several hours. The fluid may be sent for cell counts, culture, protein, or other tests. Monitoring and replacement fluids or albumin may be used according to the amount removed and the underlying illness.

\section*{Risks and aftercare}
Possible complications include pain, leakage, bleeding, infection, low blood pressure, bowel or bladder injury, and recurrence of the fluid. Keep the dressing clean and follow instructions about activity, fluids, and review. Seek advice if the drain site continues to leak or the swelling returns quickly.

\section*{When to seek urgent care}
Seek urgent help for severe or increasing abdominal pain, fever, chills, fainting, confusion, breathlessness, heavy bleeding, marked abdominal swelling, or redness and pus at the site.

\section*{Sources}
\begin{itemize}
\item Cambridge University Hospitals: \url{https://www.cuh.nhs.uk/patient-information/paracentesis/}
\end{itemize}
